% Conversão de ages-iv-sprint-0; origem imutável RELATORIO_ATUAL.odt.
\hypertarget{__RefHeading___Toc4845_3118176633}{}
\refstepcounter{subsection}
\subsection*{6.3.1 Sprint 0}
\label{ages-iv:text-h-0116}
\addcontentsline{toc}{subsection}{6.3.1 Sprint 0}

% text-p-1469
\indent A Sprint 0 marcou não apenas o início do projeto Recorda, mas também minha última participação na AGES.\originalspace{}Depois de esperar\originalspace{}muito pela chegada deste\originalspace{}momento, iniciei o projeto\originalspace{}extremamente\originalspace{}entusiasmado e com diversas ideias sobre como poderia contribuir para a organização e condução do time. Um dos primeiros desafios surgiu justamente pela composição da equipe: éramos seis AGES IV, um número elevado para uma função predominantemente gerencial. Por isso, tornou-se necessário distribuir responsabilidades de forma que pudéssemos aproveitar a experiência de todos sem criar uma estrutura excessivamente centralizada, com muitas pessoas tentando conduzir simultaneamente as mesmas decisões.\par

% text-p-1470
\indent Como uma das etapas mais importantes do início do projeto é compreender corretamente o negócio e as expectativas dos stakeholders, decidi estruturar previamente um documento reunindo dúvidas sobre a aplicação. Embora o Recorda não apresentasse uma complexidade técnica incomum, sua abrangência exigia o entendimento de diversos fluxos, funcionalidades e regras de negócio. Compartilhei o documento com o restante do time para que todos também pudessem contribuir com questionamentos e, durante a reunião com os stakeholders, nós, AGES IV, conduzimos grande parte das discussões, abordando dúvidas relacionadas tanto ao negócio quanto às decisões de tecnologia e design.\par

% text-p-1471
\indent Com apenas duas semanas para a entrega da Sprint 0, a organização inicial do time precisava acontecer de forma rápida. Optamos, então, por dividir temporariamente a equipe em grupos de trabalho. Os AGES I iniciaram as prototipações no Figma sob a liderança da Mariah, que também conduziu um workshop sobre a ferramenta. Os AGES II começaram a modelagem inicial do banco de dados, atividade que acompanhei mais diretamente, auxiliando na definição das tabelas e informações necessárias para a aplicação. Enquanto isso, os AGES III, sob a liderança do Andrei, ficaram responsáveis pelas configurações iniciais do ambiente e pelas definições da stack tecnológica, organização dos repositórios e fluxo de desenvolvimento. Os demais AGES IV concentraram-se no mapeamento dos Épicos e User Stories, buscando deixar as bases do planejamento preparadas para o início efetivo do desenvolvimento.\par

% text-p-1472
\indent Também procurei antecipar parte do planejamento das próximas Sprints. Como seria difícil reunir todos os AGES IV para discutir cada User Story antes do início das prototipações, comecei a realizar a quebra dos Épicos já definidos. Adicionei comentários e observações diretamente nas User Stories para que os demais colegas pudessem contribuir de forma assíncrona, discutindo tanto a natureza das funcionalidades quanto sua alocação ao longo das Sprints. Essa abordagem permitiu acelerar o refinamento do backlog e reduzir parte do trabalho que ainda precisaria ser realizado nas semanas seguintes.\par

% text-p-1473
\indent Uma das atividades às quais dediquei maior atenção foi a estruturação dos ambientes de gerenciamento do projeto no GitHub. Como o fluxo de desenvolvimento e versionamento também seria realizado na plataforma, aproveitamos para centralizar nela a gestão dos principais artefatos. Desenvolvi o quadro Scrumban utilizado pelo time, estruturando o fluxo das tasks entre as colunas\originalspace{}\textbf{\textmd{Product Backlog, To Do, In Progress, Ready for Review, Done e Blocked}}.\originalspace{}Além disso, foram configuradas labels relacionadas aos Épicos, Iterations de acordo com o cronograma das Sprints e filtros que facilitariam a visualização das atividades de cada ciclo. Também deixamos preparadas automações para movimentação das tasks, reduzindo parte do trabalho manual nas etapas seguintes do projeto.\par

% FIGURE draw-image-0033
\memorialfigure{imagens/original/asset-0033.png}{\textbf{Figura 33 –} Quadro Scrumban de Execução de Tasks - AGES IV}{ages-iv:draw-image-0033}{Fonte: pendente de identificação}{0.92} % <<ASSET draw-image-0033 asset-0033>>

% text-p-1476
\indent Paralelamente, construí uma visualização específica para o acompanhamento gerencial de Épicos, User Stories e Tasks. Enquanto o quadro Scrumban foi\originalspace{}pensado para facilitar a execução das atividades pelo time, essa segunda visualização oferece aos AGES IV uma perspectiva mais ampla sobre o andamento do projeto. Agrupando as informações por Épicos e User Stories, tornou-se possível identificar funcionalidades com volume desproporcional de tasks, dependências que poderiam gerar bloqueios e atividades que demandariam maior atenção ao longo das Sprints. A intenção era transformar o backlog em uma ferramenta que não apenas registrasse o planejamento, mas também ajudasse a orientar decisões gerenciais durante a execução do projeto.\par

% FIGURE draw-image-0034
\memorialfigure{imagens/original/asset-0034.png}{\textbf{Figura 34 –} Quadro Gerencial de Épicos, User Stories e Tasks\originalspace{}\originalspace{}– AGES IV}{ages-iv:draw-image-0034}{Fonte: pendente de identificação}{0.92} % <<ASSET draw-image-0034 asset-0034>>

% text-p-1479
\indent Na reta final da Sprint 0, o desenvolvimento das prototipações exigiu uma força-tarefa de praticamente todo o time. Como havia um número elevado de telas para entregar em pouco tempo, diferentes membros contribuíram diretamente para sua construção. Nesse período, os AGES II já haviam produzido uma modelagem inicial do banco de dados com o auxílio de ferramentas de IA. Para evitar que informações incorretas ou assumidas pelas ferramentas chegassem ao desenvolvimento, reservei parte do meu tempo para revisar e validar manualmente a documentação e a modelagem produzidas. Assim, conseguimos iniciar a Sprint 1 com uma base mais confiável, permitindo que os ajustes necessários fossem realizados antes que a modelagem passasse a bloquear outras atividades.\par

% text-p-1480
\indent O resultado dessa mobilização foi uma das surpresas positivas da Sprint. Mesmo diante do prazo curto, o time conseguiu prototipar todas as telas inicialmente\originalspace{}mapeadas. A proatividade demonstrada por grande parte da equipe reduziu a necessidade de um acompanhamento constante por parte dos AGES IV, que puderam concentrar seus esforços no refinamento das interfaces e na construção dos fluxos de navegação no Figma para a apresentação aos stakeholders.\par

% text-p-1481
\indent Também fiquei responsável pela elaboração e apresentação do deck da Sprint 0. Procurei construir uma apresentação objetiva e clara, apresentando os ciclos de desenvolvimento do projeto, o que havia sido realizado pela equipe e o planejamento da próxima Sprint. Minutos antes da apresentação, realizamos alguns ajustes sugeridos pela professora orientadora,\originalspace{}Alessandra Dutra,\originalspace{}incluindo um detalhamento maior das User Stories que seriam desenvolvidas. Embora inicialmente eu não enxergasse tanto valor em apresentar todas essas informações aos stakeholders, acolhi a sugestão e realizei as alterações. Ao final, recebemos um feedback positivo, especialmente em relação à clareza das informações apresentadas e à identidade visual utilizada.\par

% text-p-1482
\indent No encerramento da Sprint 0, a sensação foi de que conseguimos construir uma base sólida para o restante do projeto. Em um período curto, avançamos simultaneamente na compreensão do negócio, prototipação, modelagem inicial do banco de dados, definição da stack tecnológica, organização dos repositórios, planejamento do backlog e estruturação dos processos de gerenciamento. Mais do que entregar artefatos isolados, a Sprint serviu para organizar a forma como o Recorda seria desenvolvido ao longo dos ciclos seguintes.\par

% draw-text-box-0051
\begin{reflection}
% text-p-1484
\noindent Nesta primeira experiência como AGES IV, percebi de forma muito mais clara o peso da responsabilidade gerencial dentro de um projeto. Tenho a sensação de que o time, por melhor que seja tecnicamente, não avança sozinho: é necessário que alguém esteja constantemente pensando alguns passos à frente, preparando as próximas tasks, antecipando possíveis bloqueios e avaliando o que realmente pode ser entregue. Essa necessidade de prever cenários e criar condições para que o restante da equipe consiga trabalhar com fluidez foi algo que senti com muito mais intensidade nesta Sprint.\par

% text-p-1485
\noindent É uma responsabilidade grande, mas também uma experiência extremamente valiosa. Percebi que gerir um projeto não significa apenas distribuir tarefas ou acompanhar o que cada pessoa está fazendo, mas criar direção e preparar o caminho para que o trabalho aconteça da melhor forma possível. Acredito que essa visão será importante não apenas durante minha atuação como AGES IV, mas também na minha formação como engenheiro de software.\par

\end{reflection}

